% Fijiisland
% 2021-12-18

\frontmatter
\chapter{序言}

Rust是一门为系统编程而生的语言。
\mkpar{1}
然而在如今，绝大多数正在工作的程序员对于系统编程并不熟悉，那么便需要对其进行一些解释了。至今它仍是我们做任何事的基础。
\mkpar{1}
你合上了笔记本电脑。操作系统会感知到，并中止所有正在运行的程序，熄灭屏幕，最后让电脑进入睡眠。过了一会儿，你又打开了
电脑，屏幕与其它组件再次被启动，并且每个程序均能从它中断的地方恢复过来，我们认为这是理所当然的事。但这都是系统程序员
写了大量代码而得到的成果。
\mkpar{1}
系统编程一般被用到这些地方：
\begin{itemize}
    \item 操作系统
    \item 各种设备驱动
    \item 文件系统
    \item 数据库
    \item 运行在廉价设备、亦或是极度可靠的设备上的代码
    \item 密码学
    \item 媒体解码器（负责读取及写入音频、视频或者图像文件的软件）
    \item 媒体处理（例如，语音识别或照片编辑软件）
    \item 内存管理（例如，实现一个垃圾回收器）
    \item 文本渲染（将文本和字体转变为像素）
    \item 实现高级编程语言（比如 JavaScript 和 Python）
    \item 网络
    \item 虚拟化技术以及软件容器
    \item 科学模拟
    \item 游戏
\end{itemize}
简单来说，系统编程是一种资源受限的程序设计，它面向每个字节（byte）和每个CPU周期（cycle）来编程。
\mkpar{1}
涉及支持一个基础应用程序运行的系统代码的数量会相当惊人。

\mkbaresec{适合阅读本书的人群}
如果你已是一位系统程序员并且准备好迎接C++的替代品，那么本书适合你。若你是对于任何编程语言均有经验的开发者，无论是
C\#，Java，Python，JavaScript或者是其它东西，那么本书也适合你。
\mkpar{1}
但是，你不能光学习Rust。要想最大化地利用该门语言，你还需获得一些系统编程的经验。我们建议你在阅读本书的同时用Rust
实现一些有关系统编程的项目。试着去做一些你从来没有构建过的，能利用好Rust的迅捷、并发以及安全性的东西。本序言开端列
出的这些议题应该能为你带来一些灵感。

\mkbaresec{我们编写本书的原因}
我们想编写一本我们自己当初学习Rust时便希望拥有的书。我们的目标是开门见山地，将Rust重要且新的概念整理出来，进而清晰
而深入地表述它们，以便减少“通过试错来学习”这样的现象\textbf{（待修订）}。

\mkbaresec{本书导航}
本书的前两章介绍了Rust，并在我们移步至\boldred{第三章}的基础数据类型前提供了一段简短的总览。有关所有权和引用的概念
位于第\boldred{四、五}章节。我们建议按顺序通读本书前五章。
\mkpar{1}
第\boldred{六}至第\boldred{十}章涵盖了该门语言的基础知识：表达式（\boldred{第六章}）、错误处理（\boldred{第七章}）、程序包
和模块（\boldred{第八章}）、结构体（\boldred{第九章}）以及枚举和模式（\boldred{第十章}）。对这些章进行跳跃式阅读没问题，但一
定别跳过错误处理这章，相信我们。
\mkpar{1}
\boldred{第十一章}涵盖了特型和泛型，也是最后两个你需要了解的重要概念。特型类似于Java或C\#里的接口。同时它们作为一种主要方式，使得
Rust能够允许你将你的自定义类型整合进该语言本身。\boldred{第十二章}展示了特型是如何支持运算符重载的，另外\boldred{第十三章}涵盖了
更多实用特型。
\mkpar{1}
了解特型和泛型之后便可解锁本书的余下部分。第\boldred{十四、十五}章涵盖了闭包和迭代器这两个你绝对
不想错过的强大工具。对于剩余章节你可以按自由顺序进行阅读，或者根据需求对它们进行深度挖掘。这些章节涵盖了该门语言的剩余部分：集合（\boldred{第十六章}）、字符串和文本（\boldred{第十七章}）、输
入输出（\boldred{第十八章}）、并发（\boldred{第十九章}）、异步编码（\boldred{第二十章}）、宏（\boldred{第二十一章}）、unsafe编码（\boldred{第二十二章}）以及如何在其它语言中调用Rust函数（\boldred{第二十三章}）。

\mkbaresec{本书中的惯例用法}
以下排版惯例被用于本书中：
\mkpar{1}
意大利斜体 --- \textit{Italic}
    \begin{quote}
        表示新术语、网址、电子邮件地址、文件名以及文件拓展名。
    \end{quote}
\mkpar{1}
等宽体 --- \texttt{Constant width}
    \begin{quote}
        用来列出程序代码，也用作在段内引用如变量、函数名、数据库、数据类型、环境变量、声明式以及关键字这样的程序元素。
    \end{quote}
\mkpar{1}
等宽粗体 --- \texttt{\textbf{Constant width bold}}
\begin{quote}
    表示需由用户键入的命令或其它文本。
\end{quote}
\mkpar{1}
等宽斜体 --- \textit{\texttt{Constant width italic}}
\begin{quote}
    表示需要被替换的文本，由来自用户的值或由上下文确定的值提供。
\end{quote}
\mkpar{1}
\mknote{该框体代表一个一般注释。}

\mkbaresec{使用示例代码}
补充材料（示例代码、编程练习等）可于下方网址下载
\mkpar{0.5}
\refred{https://github.com/ProgrammingRust.}
\mkpar{0.5}
本书是为了帮你完成工作的。通常来说，若示例代码由本书提供，你可以将它用于你的项目和文档中。若非大量地复制本书中的代码，那么
你不需要联系我们以获取授权。比如说，
